\subsection{Linear Projective Geometry and Affine Completion}

Throughout this subsection, $\mathcal{V}$ is finite-dimensional over a fixed field $\mathbb{F}$.

\begin{definition} \label{def:projective-geometry-vector-space} \hlt{Projective Geometry of a Vector Space}
\begin{enumerate}[label=(\roman*)]
\setlength{\itemsep}{0pt}
\item The \hlt{projective geometry} of $\mathcal{V}$ is the poset $\operatorname{Sub}(\mathcal{V})$ of linear subspaces, ordered by inclusion.\\
Its elements are \hlt{projective subspaces}, and its dimension map is $\operatorname{pdim}:\operatorname{Sub}(\mathcal{V})\to\mathbb{Z}$, $\operatorname{pdim}(\mathcal{U})\defeq\dim(\mathcal{U})-1$.
\item A \hlt{projective point} is a $1$-dimensional subspace; a \hlt{projective line} is a $2$-dimensional subspace.
\end{enumerate}
\end{definition}

\begin{theorem} \label{thm:linear-span-map-affine-to-projective} \hlt{Linear Span Map from Affine to Projective Geometry}\\
Let $H=\mathbf{x}+\mathcal{K}$ be an affine hyperplane with $\mathbf{0}\notin H$, and put $\mathcal{A}(H)\defeq\{X\in\mathcal{A}(\mathcal{V}):X\subseteq H\}$.\\
Define $P:\mathcal{A}(H)\to\operatorname{Sub}(\mathcal{V})$ by $P(X)\defeq\operatorname{span}(X)$. Then:
\begin{enumerate}[label=(\roman*)]
\setlength{\itemsep}{0pt}
\item $P$ is injective, and $(\forall\,\mathcal{U}\in\operatorname{im}P)\,P^{-1}(\mathcal{U})=\mathcal{U}\cap H$.
\item $\operatorname{im}P=\{\mathcal{U}\leq\mathcal{V}:\mathcal{U}\nsubseteq\mathcal{K}\}$.
\item $(\forall\,X,Y\in\mathcal{A}(H))\,[X\subseteq Y\iff P(X)\subseteq P(Y)]$.
\item $(\forall\,\mathcal{C}\subseteq\mathcal{A}(H))\,[\mathcal{C}\neq\emptyset\land\bigcap\mathcal{C}\neq\emptyset\Rightarrow P(\bigcap\mathcal{C})=\bigcap\limits_{X\in\mathcal{C}}P(X)]$.
\item $(\forall\,\mathcal{C}\subseteq\mathcal{A}(H))\,[\mathcal{C}\neq\emptyset\Rightarrow P(\bigvee\mathcal{C})=\bigvee\limits_{X\in\mathcal{C}}P(X)=\sum\limits_{X\in\mathcal{C}}P(X)]$, where $\bigvee\mathcal{C}$ is computed in $\mathcal{A}_0(\mathcal{V})$.
\item $(\forall\,X\in\mathcal{A}(H))\,\operatorname{pdim}(P(X))=\dim(X)$.
\item $(\forall\,X,Y\in\mathcal{A}(H))\,[X\parallel Y\iff(P(X)\cap\mathcal{K}\subseteq P(Y)\cap\mathcal{K}\lor P(Y)\cap\mathcal{K}\subseteq P(X)\cap\mathcal{K})]$.
\end{enumerate}
\begin{proof}
\begin{enumerate}[label=(\roman*)]
\setlength{\itemsep}{0pt}
\item[$(i)$] Fix $X=\mathbf{u}+\mathcal{S}\in\mathcal{A}(H)$. Then $\mathcal{S}\leq\mathcal{K}$ by Proposition \ref{prop:basic-flat-incidence}(iii).\\
If $\mathbf{u}\in\mathcal{K}$, then $\mathbf{u}-\mathbf{x}\in\mathcal{K}$ because $\mathbf{u}\in H$, so $\mathbf{x}\in\mathcal{K}$ and $\mathbf{0}\in H$, a contradiction.\\
Thus $\mathbf{u}\notin\mathcal{K}$. Since $\mathbf{u}\in X$ and $(\forall\,\mathbf{s}\in\mathcal{S})\,\mathbf{s}=(\mathbf{u}+\mathbf{s})-\mathbf{u}$, $\{\mathbf{u}\}\cup\mathcal{S}\subseteq P(X)$.\\
Together with $X\subseteq\mathbb{F}\mathbf{u}+\mathcal{S}$, this proves $P(X)=\mathbb{F}\mathbf{u}\oplus\mathcal{S}$; the sum is direct because $\mathcal{S}\leq\mathcal{K}$ and $\mathbf{u}\notin\mathcal{K}$.\\
If $\mathbf{z}\in P(X)\cap H$, write $\mathbf{z}=r\mathbf{u}+\mathbf{s}=\mathbf{u}+\mathbf{k}$ with $(r,\mathbf{s},\mathbf{k})\in\mathbb{F}\times\mathcal{S}\times\mathcal{K}$.\\
Then $(1-r)\mathbf{u}=\mathbf{s}-\mathbf{k}\in\mathcal{K}$, whence $r=1$ and $\mathbf{z}\in X$. Since $X\subseteq P(X)\cap H$, $P(X)\cap H=X$.\\
Therefore $(\forall\,X,Y\in\mathcal{A}(H))\,[P(X)=P(Y)\Rightarrow X=P(X)\cap H=P(Y)\cap H=Y]$, and $(\forall\,\mathcal{U}\in\operatorname{im}P)\,P^{-1}(\mathcal{U})=\mathcal{U}\cap H$.
\item[$(ii)$] By the representation in (i), $(\forall\,X\in\mathcal{A}(H))\,P(X)\nsubseteq\mathcal{K}$. Conversely, fix $\mathcal{U}\leq\mathcal{V}$ with $\mathcal{U}\nsubseteq\mathcal{K}$ and choose $\mathbf{u}\in\mathcal{U}\setminus\mathcal{K}$.\\
Since $\mathcal{K}$ has codimension one, $\dim(\mathcal{V}/\mathcal{K})=1$ by Theorem \ref{thm:dimension-of-quotient}(ii).\\
The cosets $\mathbf{x}+\mathcal{K}$ and $\mathbf{u}+\mathcal{K}$ are nonzero; choose $r\in\mathbb{F}^{\times}$ with $\mathbf{x}+\mathcal{K}=r(\mathbf{u}+\mathcal{K})$, and put $\mathbf{h}\defeq r\mathbf{u}\in\mathcal{U}\cap H$.\\
Then $(\forall\,\mathbf{z}\in\mathcal{V})\,[\mathbf{z}\in\mathcal{U}\cap H\iff\mathbf{z}-\mathbf{h}\in\mathcal{U}\cap\mathcal{K}]$, so $\mathcal{U}\cap H=\mathbf{h}+(\mathcal{U}\cap\mathcal{K})\in\mathcal{A}(H)$.\\
By $\operatorname{span}(\mathbf{h}+\mathcal{K})=\mathcal{V}/\mathcal{K}$, $(\forall\,\mathbf{w}\in\mathcal{U})(\exists\,a\in\mathbb{F})\,\mathbf{w}-a\mathbf{h}\in\mathcal{U}\cap\mathcal{K}$. Therefore $P(\mathcal{U}\cap H)=\mathbb{F}\mathbf{h}+(\mathcal{U}\cap\mathcal{K})=\mathcal{U}$.
\item[$(iii)$] ($\Rightarrow$) Fix $X,Y\in\mathcal{A}(H)$. Then $X\subseteq Y\Rightarrow\operatorname{span}(X)\subseteq\operatorname{span}(Y)$.
\item[$(iii)$] ($\Leftarrow$) Fix $X,Y\in\mathcal{A}(H)$. By (i), $P(X)\subseteq P(Y)\Rightarrow X=P(X)\cap H\subseteq P(Y)\cap H=Y$.
\item[$(iv)$] Fix $\emptyset\neq\mathcal{C}\subseteq\mathcal{A}(H)$ with $\bigcap\mathcal{C}\neq\emptyset$, and choose $\mathbf{z}\in\bigcap\mathcal{C}$. For $X\in\mathcal{C}$, let $\mathcal{S}_X$ be its base; then $X=\mathbf{z}+\mathcal{S}_X$ by Proposition \ref{prop:flat-base-and-representatives}(ii).\\
Hence $\bigcap\mathcal{C}=\mathbf{z}+\bigcap\limits_{X\in\mathcal{C}}\mathcal{S}_X$ by Lemma \ref{lem:intersection-of-flats}. Since $\mathbf{z}\notin\mathcal{K}$, $(\forall\,X\in\mathcal{C})\,\mathcal{S}_X\leq\mathcal{K}$, $\bigcap\limits_{X\in\mathcal{C}}P(X)=\bigcap\limits_{X\in\mathcal{C}}(\mathbb{F}\mathbf{z}\oplus\mathcal{S}_X)=\mathbb{F}\mathbf{z}\oplus\bigcap\limits_{X\in\mathcal{C}}\mathcal{S}_X=P(\bigcap\mathcal{C})$.
\item[$(v)$] Fix $\emptyset\neq\mathcal{C}\subseteq\mathcal{A}(H)$, put $E\defeq\bigcup\mathcal{C}$. By Theorem \ref{thm:lattice-of-flats}(ii), $\bigvee\mathcal{C}=\operatorname{affhull}(E)\subseteq H$ as $H$ is upper bound, so $\bigvee\mathcal{C}\in\mathcal{A}(H)$.\\
By Theorem \ref{thm:affine-hull-span-formula}(i), $\operatorname{affhull}(E)=\operatorname{affspan}(E)\subseteq\operatorname{span}(E)$. By $E\subseteq\operatorname{affhull}(E)$, $\operatorname{span}(E)\subseteq\operatorname{span}(\operatorname{affhull}(E))$.\\
Hence $P(\bigvee\mathcal{C})=\operatorname{span}(E)=\sum\limits_{X\in\mathcal{C}}\operatorname{span}(X)=\sum\limits_{X\in\mathcal{C}}P(X)=\bigvee\limits_{X\in\mathcal{C}}P(X)$.
\item[$(vi)$] If $X=\mathbf{u}+\mathcal{S}\in\mathcal{A}(H)$, then $P(X)=\mathbb{F}\mathbf{u}\oplus\mathcal{S}$ by (i), so $\operatorname{pdim}(P(X))=\dim(\mathbb{F}\mathbf{u}\oplus\mathcal{S})-1=\dim\mathcal{S}=\dim X$.
\item[$(vii)$] ($\Rightarrow$) Fix $X=\mathbf{u}+\mathcal{S},Y=\mathbf{v}+\mathcal{T}\in\mathcal{A}(H)$. By the calculation in (i), $P(X)\cap\mathcal{K}=\mathcal{S}$ and $P(Y)\cap\mathcal{K}=\mathcal{T}$.\\
By $X\parallel Y$ and the definition of parallel flats, $\mathcal{S}\subseteq\mathcal{T}$ or $\mathcal{T}\subseteq\mathcal{S}$, so the two intersections with $\mathcal{K}$ are comparable.
\item[$(vii)$] ($\Leftarrow$) Fix $X=\mathbf{u}+\mathcal{S},Y=\mathbf{v}+\mathcal{T}\in\mathcal{A}(H)$. By the identities in the forward direction, comparability of $P(X)\cap\mathcal{K}$ and $P(Y)\cap\mathcal{K}$ is comparability of $\mathcal{S}$ and $\mathcal{T}$. Hence $X\parallel Y$ by the definition of parallel flats.
\end{enumerate}
\hlt{Structural Commentary}: Compute spans and traces on $H$, then transport order, lattice operations, dimension, and direction.\\
\hlt{Key Insight}: Transversality makes intersection with $H$ inverse to taking linear span.
\end{proof}
\end{theorem}

\begin{perspective} \label{per:projective-completion-from-affine-hyperplane} \hlt{Projective Completion of an Affine Hyperplane}\\
Let $H=\mathbf{x}+\mathcal{K}$ be an affine hyperplane with $\mathbf{0}\notin H$. Theorem \ref{thm:linear-span-map-affine-to-projective} identifies each affine flat $X\subseteq H$ with $\operatorname{span}(X)$ and records its direction as $\operatorname{span}(X)\cap\mathcal{K}$. If $\dim H=2$, distinct parallel affine lines therefore determine the same one-dimensional subspace of $\mathcal{K}$, their common projective point at infinity. This viewpoint explains the completion; the theorem owns the incidence claims.
\end{perspective}

Figure \ref{fig:affine-projective-span-hyperplane} illustrates the preceding plane case: parallel affine lines are disjoint in $H$, whereas their projective spans meet in a shared direction contained in $\mathcal{K}$.

\begin{illustrationfigure}[H]
  \centering
  % Schematic: affine flats in H=x+K become linear/projective flats by taking spans.
\begin{tikzpicture}[
  x=0.88cm,
  y=0.86cm,
  >=Stealth,
  every node/.style={font=\small},
  plane/.style={draw=gray!62!black, fill=gray!9, line width=0.55pt},
  spanx/.style={draw=blue!58!black, fill=blue!12, fill opacity=0.28, line width=0.55pt},
  spany/.style={draw=teal!58!black, fill=teal!10, fill opacity=0.24, line width=0.55pt},
  affine/.style={blue!72!black, line width=1.1pt},
  direction/.style={teal!65!black, dashed, line width=0.9pt},
  guide/.style={gray!55, line width=0.45pt},
  point/.style={circle, fill=gray!75!black, inner sep=1.35pt},
  labelnode/.style={font=\scriptsize, fill=white, fill opacity=1, text opacity=1, inner sep=1.35pt}
]
  \coordinate (K1) at (-3.15,-1.55);
  \coordinate (K2) at (2.85,-1.55);
  \coordinate (K3) at (3.48,-1.05);
  \coordinate (K4) at (-2.52,-1.05);

  \coordinate (H1) at (-3.15,0.85);
  \coordinate (H2) at (2.85,0.85);
  \coordinate (H3) at (3.48,1.35);
  \coordinate (H4) at (-2.52,1.35);

  \coordinate (D1) at (-1.85,-1.30);
  \coordinate (D2) at (1.65,-1.30);
  \coordinate (O) at (-0.10,-1.30);

  \coordinate (X1) at (-2.05,1.02);
  \coordinate (X2) at (1.45,1.02);
  \coordinate (Y1) at (-1.15,1.20);
  \coordinate (Y2) at (2.35,1.20);

  \filldraw[plane] (K1) -- (K2) -- (K3) -- (K4) -- cycle;
  \node[labelnode] at (3.12,-1.10) {$\mathcal K$};

  \filldraw[plane] (H1) -- (H2) -- (H3) -- (H4) -- cycle;
  \node[labelnode] at (2.55,1.55) {$H=\mathbf{x}+\mathcal K$};

  \filldraw[spanx] (D1) -- (D2) -- (X2) -- (X1) -- cycle;
  \filldraw[spany] (D1) -- (D2) -- (Y2) -- (Y1) -- cycle;

  \draw[guide] (D1) -- (X1);
  \draw[guide] (D2) -- (X2);
  \draw[guide] (D1) -- (Y1);
  \draw[guide] (D2) -- (Y2);

  \draw[direction] (D1) -- (D2);
  \node[labelnode, text=teal!60!black, align=center] at (0.70,-1.08)
    {$D=\operatorname{span}(X)\cap\operatorname{span}(Y)$};

  \draw[affine] (X1) -- (X2);
  \draw[affine] (Y1) -- (Y2);
  \node[labelnode, text=blue!70!black] at (-2.25,0.78) {$X$};
  \node[labelnode, text=blue!70!black] at (2.45,1.03) {$Y$};
  \node[labelnode, text=blue!58!black] at (-0.45,0.00) {$\operatorname{span}(X)$};
  \node[labelnode, text=teal!58!black] at (0.95,0.36) {$\operatorname{span}(Y)$};

  \node[point] at (O) {};
  \node[labelnode] at (-0.25,-1.58) {$\mathbf 0$};
\end{tikzpicture}

  \caption{The affine plane $H=\mathbf{x}+\mathcal{K}$, direction plane $\mathcal{K}$, parallel lines $X,Y$, and shaded spans $\operatorname{span}(X),\operatorname{span}(Y)$. The spans meet along the dashed subspace $D=\operatorname{span}(X)\cap\mathcal{K}=\operatorname{span}(Y)\cap\mathcal{K}$, the common projective point at infinity.}
  \label{fig:affine-projective-span-hyperplane}
\end{illustrationfigure}
